% !TEX root =./ECE448_Practice_main.tex
%
% L14 practice set (HD Radio). Every question is a step in finding, measuring,
% or decoding the hybrid sidebands: locating subcarriers on a real station,
% choosing an FFT size, the OFDM timing and block parameters, and reading sideband
% power. The L14 Lab asks for the extended hybrid edges, the FFT sizes 1024 and
% 8192 at 1,488,375 samples per second, and the -14 dBc case, so none of
% those appear here.
% Blank copy: \begin{solution}[h] reserves h of work space.
% SOLUTIONS copy (\iskey=1) prints the red worked solutions.

\fullwidth{\textbf{Learning Objective 6.3: } I can describe how HD Radio works and
demodulate an HD Radio signal.
\\[4pt]\normalfont\small Show your work. Write a \textbf{Documentation} line at the
end (\textbf{None} if you did not collaborate). Use the hybrid numbers from the lesson:
subcarrier spacing $\Delta f = 1{,}488{,}375/4096$ Hz, sidebands from subcarrier 356 to 546
on each side, ten partitions of 19 subcarriers per sideband, and a guard of $7/128$ of the
useful symbol time.}

\begin{questions}

%%%%%%%%%%%%%%%%%%%%%%%%%%%%%%%%%%%%%%%%%
\question \textbf{Finding the sidebands.} A hybrid station broadcasts at 98.5 MHz.
\begin{parts}

	\begin{minipage}[t]{0.58\linewidth}
		\part Find the frequency of subcarriers 356, 450, and 546, in kHz from the channel
		center, to 0.1 kHz.
		\begin{solution}[0.9in]
		\eq{ \Delta f &= 1{,}488{,}375/4096 = 363.37\ \text{Hz} \\
			356\,\Delta f &= 129.4\ \text{kHz} \\
			450\,\Delta f &= 163.5\ \text{kHz} \\
			546\,\Delta f &= 198.4\ \text{kHz} }
		\end{solution}
	\end{minipage}\hfill%
	\begin{minipage}[t]{0.37\linewidth}
		\raggedleft \vspace{12pt}
		\ansbox[1.6in]{129.4, 163.5, 198.4 kHz}
	\end{minipage}

	\begin{minipage}[t]{0.58\linewidth}
		\part Between which two radio frequencies, in MHz, should the upper sideband appear on
		a spectrum display? The lower sideband?
		\begin{solution}[0.8in]
		Upper: $98.5 + 0.1294 = 98.6294$ MHz to $98.5 + 0.1984 = 98.6984$ MHz. Lower:
		98.3016 to 98.3706 MHz.
		\end{solution}
	\end{minipage}\hfill%
	\begin{minipage}[t]{0.37\linewidth}
		\raggedleft \vspace{12pt}
		upper \ansbox[1.5in]{98.6294 to 98.6984}

		\vspace{4pt}
		lower \ansbox[1.5in]{98.3016 to 98.3706}
	\end{minipage}

	\part Counting the partition that starts at subcarrier 356 as the first, which partition
	holds subcarrier 450? The reference subcarriers are 356, 375, 394, and every 19th after
	that, up to 546. Is subcarrier 451 a reference subcarrier?
		\begin{solution}[0.9in]
		Partition $p$ covers subcarriers $356 + 19(p-1)$ to $356 + 19p - 1$. Since
		$450 - 356 = 94 = 4 \times 19 + 18$, subcarrier 450 is the last subcarrier of the fifth
		partition, 432 to 450. Since $451 - 356 = 95 = 5 \times 19$, subcarrier 451 is a
		reference subcarrier, the first subcarrier of the sixth partition.
		\end{solution}

\end{parts}

\newpage
%%%%%%%%%%%%%%%%%%%%%%%%%%%%%%%%%%%%%%%%%
\question \textbf{Choosing an FFT size.} You record the band with an RTL-SDR at
2.4 MS/s and look at it with a QT GUI Frequency Sink.
\begin{parts}

	\begin{minipage}[t]{0.58\linewidth}
		\part Find the bin width and the number of bins one 69,041 Hz sideband spans, for an
		FFT of 2048 points and of 16384 points.
		\begin{solution}[1.1in]
		\eq{ 2048:&\ 2.4\ \text{MHz}/2048 = 1171.9\ \text{Hz} \\
			&\ 69{,}041/1171.9 = 58.9\ \text{bins} \\
			16384:&\ 146.5\ \text{Hz},\ 471\ \text{bins} }
		\end{solution}
	\end{minipage}\hfill%
	\begin{minipage}[t]{0.37\linewidth}
		\raggedleft \vspace{12pt}
		2048 \ansbox[1.3in]{1171.9 Hz, 59}

		\vspace{4pt}
		16384 \ansbox[1.3in]{146.5 Hz, 471}
	\end{minipage}

	\part Which FFT size can locate a sideband's edge to within one 363.4 Hz subcarrier?
	Which one would you use to find the sidebands quickly on a live band, and why? Why does
	neither show the subcarriers as separate peaks?
		\begin{solution}[1.3in]
		Only 16384: its 146.5 Hz bins are narrower than a subcarrier, while a 2048-point bin is
		3.2 subcarriers wide. To find the sidebands quickly, the 2048-point FFT is enough: 59
		bins still draw a clear flat shelf, the display updates and averages faster, and the
		edges are already within 1.2 kHz. Neither shows separate peaks because the subcarriers
		are orthogonal: their sinc spectra overlap and add to a flat top, with no dip between
		neighbors.
		\end{solution}

\end{parts}

\newpage
%%%%%%%%%%%%%%%%%%%%%%%%%%%%%%%%%%%%%%%%%
\question \textbf{OFDM timing and block parameters.}
\begin{parts}

	\begin{minipage}[t]{0.58\linewidth}
		\part Find the useful symbol time $T_u$, the guard $T_g$, the whole symbol time, and
		the symbol rate.
		\begin{solution}[1.2in]
		\eq{ T_u &= 1/\Delta f = 2.752\ \text{ms} \\
			T_g &= (7/128)\,T_u = 150.5\ \mu\text{s} \\
			T_s &= T_u + T_g = 2.902\ \text{ms} \\
			R_s &= 1/T_s = 344.5\ \text{symbols/s} }
		\end{solution}
	\end{minipage}\hfill%
	\begin{minipage}[t]{0.37\linewidth}
		\raggedleft \vspace{12pt}
		$T_u$ \ansbox[1.2in]{2.752 ms}

		\vspace{4pt}
		$T_g$ \ansbox[1.2in]{150.5 $\mu$s}

		\vspace{4pt}
		$R_s$ \ansbox[1.2in]{344.5 /s}
	\end{minipage}

	\begin{minipage}[t]{0.58\linewidth}
		\part nrsc5 reads complex float recordings at 744,187.5 samples per second. In a GNU
		Radio OFDM transmitter at that rate, what \texttt{fft\_len} puts one subcarrier in each
		FFT bin, and what \texttt{cp\_len} gives NRSC-5's guard of $7/128$ of the FFT?
		\begin{solution}[0.9in]
		\eq{ \texttt{fft\_len} &= 744{,}187.5/363.37 = 2048 \\
			\texttt{cp\_len} &= (7/128) \times 2048 = 112 }
		\end{solution}
	\end{minipage}\hfill%
	\begin{minipage}[t]{0.37\linewidth}
		\raggedleft \vspace{12pt}
		\texttt{fft\_len} \ansbox[1.1in]{2048}

		\vspace{4pt}
		\texttt{cp\_len} \ansbox[1.1in]{112}
	\end{minipage}

	\begin{minipage}[t]{0.58\linewidth}
		\part At what sample rate would \texttt{fft\_len} = 8192 put one subcarrier in each
		bin, and what \texttt{cp\_len} goes with it? Could the course's RTL-SDRs, which run at
		2.4 MS/s, record at that rate?
		\begin{solution}[0.9in]
		\eq{ 8192 \times 363.37 &= 2{,}976{,}750\ \text{S/s} \\
			\texttt{cp\_len} &= (7/128) \times 8192 = 448 }
		No: that is above the 2.4 MS/s the course uses, so the RTL-SDR records at
		1,488,375 S/s with \texttt{fft\_len} = 4096 instead.
		\end{solution}
	\end{minipage}\hfill%
	\begin{minipage}[t]{0.37\linewidth}
		\raggedleft \vspace{12pt}
		rate \ansbox[1.2in]{2,976,750 S/s}

		\vspace{4pt}
		\texttt{cp\_len} \ansbox[1.2in]{448}
	\end{minipage}

	\begin{minipage}[t]{0.58\linewidth}
		\part A reflection from a ridge travels 18 km farther than the direct path. Find its
		delay. Does it fit inside the guard, whose edges NRSC-5 tapers? How many symbols would
		it span if the sideband were one carrier with a 14.5 $\mu$s symbol?
		\begin{solution}[1.0in]
		\eq{ \tau = 18\ \text{km}/(3\times10^8\ \text{m/s}) = 60\ \mu\text{s} }
		It fits in the first half of the 150.5 $\mu$s guard, clear of the tapered edge. On
		one carrier it would span $60/14.5 \approx 4$ symbols.
		\end{solution}
	\end{minipage}\hfill%
	\begin{minipage}[t]{0.37\linewidth}
		\raggedleft \vspace{12pt}
		$\tau$ \ansbox[1.1in]{60 $\mu$s}

		\vspace{4pt}
		symbols \ansbox[1.1in]{about 4}
	\end{minipage}

\end{parts}

\newpage
%%%%%%%%%%%%%%%%%%%%%%%%%%%%%%%%%%%%%%%%%
\question \textbf{Sideband power.} An FM station's analog signal is 6 kW effective
radiated power.
\begin{parts}

	\begin{minipage}[t]{0.58\linewidth}
		\part At $-20$ dBc, find the total digital power in watts, the power of each sideband
		in watts and in dBc, and the power of each of the 382 subcarriers in dBc.
		\begin{solution}[1.1in]
		\eq{ 6000 \times 10^{-2} &= 60\ \text{W total, 30 W per sideband} \\
			-20 - 10\log_{10}2 &= -23.0\ \text{dBc} \\
			-20 - 10\log_{10}382 &= -45.8\ \text{dBc} }
		\end{solution}
	\end{minipage}\hfill%
	\begin{minipage}[t]{0.37\linewidth}
		\raggedleft \vspace{12pt}
		total \ansbox[1.2in]{60 W}

		\vspace{4pt}
		sideband \ansbox[1.2in]{30 W, $-23.0$ dBc}

		\vspace{4pt}
		subcarrier \ansbox[1.2in]{$-45.8$ dBc}
	\end{minipage}

	\begin{minipage}[t]{0.58\linewidth}
		\part The station raises its digital power to $-10$ dBc total. Repeat (a).
		\begin{solution}[0.9in]
		600 W total and 300 W per sideband, which is $-13.0$ dBc each and $-35.8$ dBc per
		subcarrier.
		\end{solution}
	\end{minipage}\hfill%
	\begin{minipage}[t]{0.37\linewidth}
		\raggedleft \vspace{12pt}
		\ansbox[1.6in]{600 W, $-13.0$, $-35.8$ dBc}
	\end{minipage}

	\begin{minipage}[t]{0.58\linewidth}
		\part The L14 meters read $-26.5$ dBc for the upper sideband and $-23.1$ dBc for the
		lower on another station. Find the total digital power in dBc, and say what the
		difference between the two tells you.
		\begin{solution}[1.0in]
		\eq{ 10\log_{10}\left(10^{-2.65} + 10^{-2.31}\right) = -21.5\ \text{dBc} }
		The sidebands differ by 3.4 dB, so the station runs asymmetric sidebands, most likely
		to protect a neighbor above it.
		\end{solution}
	\end{minipage}\hfill%
	\begin{minipage}[t]{0.37\linewidth}
		\raggedleft \vspace{12pt}
		\ansbox[1.2in]{$-21.5$ dBc}
	\end{minipage}

	\part On a weak station the shelves stand only 10 dB above the noise floor. Does the
	meter read the sideband high or low, and by about how much? Why does the frequency sink
	show the shelves about 25 dB under the host's peak when each sideband is $-23$ dBc?
		\begin{solution}[1.3in]
		High: the meter adds the noise in the band to the sideband's power, and noise 10 dB
		down adds 10\%, which is
		\eq{ 10\log_{10}1.1 = 0.41\ \text{dB} }
		The frequency sink shows power per bin, not power per band. The host's power is
		concentrated in the bins near its center, while each sideband spreads its power evenly
		over 191 bins, so the host's peak bin is higher relative to a sideband bin than the
		host's total power is relative to the sideband's total.
		\end{solution}

\end{parts}

\end{questions}

